\section{Step 1 Predict a rectangular transpose}\label{transpose-s01}

Write the three input rows on paper before editing a loop. The dimensions belong to the representation: the same fifteen words laid out as five rows of three describe the output, not the input. A source coordinate and an output coordinate name the same value through different row widths.

\subsection{Write or inspect}
Read source row one, column zero in the worked example. Locate its destination in the five-row output. Then predict the entire first output row.

\subsection{Expected observation}
The value five moves to output row zero, column one. The first output row is zero, five, ten.

\Needspace{11\baselineskip}
\subsection{Check your understanding}
Which output coordinate receives input coordinate (1, 0)?

\begin{itemize}
\item \textbf{A} (0, 1).
\item \textbf{B} (1, 0).
\item \textbf{C} (1, 4).
\end{itemize}
